\section{ANEXOS Y GUÍA DEL PROCEDIMIENTO TÉCNICO DE REPLICABILIDAD}

En esta sección se detalla el procedimiento técnico paso a paso para la replicabilidad del experimento, la configuración del entorno de ejecución, las especificaciones de hardware y las instrucciones de acceso al código fuente hospedado en el repositorio oficial.

\subsection{Enlace Oficial al Repositorio de Código Fuente en GitHub}
La totalidad de los códigos fuente, scripts de procesamiento, modelos entrenados, gráficos en alta resolución y cuadernos de Jupyter compilados se encuentran alojados en el repositorio público de GitHub:
\begin{center}
\textbf{URL del Repositorio de GitHub}: \url{https://github.com/carmenNieves6478/EdgeGuard--IoT.git}
\end{center}

\subsection{Anexo A: Procedimiento Técnico de Ingesta y Limpieza del Dataset}
\begin{enumerate}[leftmargin=*]
    \item \textbf{Ubicación de Datos Brutos}: Los 63 archivos CSV del dataset CIC-IoT-2023 se ubican en la ruta local \texttt{E:\textbackslash PROYECTO DE INVESTIGACION\textbackslash dataset\textbackslash MERGED\_CSV\textbackslash MERGED\_CSV}.
    \item \textbf{Ejecución del Script de Preprocesamiento}: Se ejecuta el script de ingesta por fragmentos:
    \begin{lstlisting}[language=bash]
    python process_dataset.py
    \end{lstlisting}
    \item \textbf{Generación de Artefactos}: El procedimiento produce los archivos Parquet optimizados \texttt{train.parquet} ($130,424$ filas), \texttt{val.parquet} ($27,948$ filas) y \texttt{test.parquet} ($27,949$ filas), así como los objetos de escalado \texttt{scaler.joblib} y codificación \texttt{label\_encoder.joblib} en la carpeta \texttt{models/}.
\end{enumerate}

\subsection{Anexo B: Procedimiento de Entrenamiento del Modelo VAE-MLP en GPU}
\begin{enumerate}[leftmargin=*]
    \item \textbf{Ejecución del Entrenamiento PyTorch}: Se ejecuta el script de entrenamiento en GPU CUDA durante 40 épocas:
    \begin{lstlisting}[language=bash]
    python train_vae_mlp.py
    \end{lstlisting}
    \item \textbf{Exportación y Cuantificación ONNX INT8}: Se ejecuta la exportación al runtime ONNX y cuantificación dinámicas pos-entrenamiento:
    \begin{lstlisting}[language=bash]
    python quantize_and_xai.py
    \end{lstlisting}
    \item \textbf{Verificación de Artefactos de Salida}: Produce el modelo PyTorch \texttt{vae\_mlp\_best.pt}, el modelo ONNX cuantizado \texttt{vae\_mlp\_quantized.onnx} ($21.67\text{ KB}$) y los gráficos de importancia SHAP.
\end{enumerate}

\subsection{Anexo C: Procedimiento de Entrenamiento del Ensamble Stacking}
\begin{enumerate}[leftmargin=*]
    \item \textbf{Entrenamiento de Clasificadores Nivel 0 y Meta-Learner LightGBM}: Se ejecuta el script de optimización de artefactos de ensamble:
    \begin{lstlisting}[language=bash]
    python save_optimized_artifacts.py
    \end{lstlisting}
    \item \textbf{Salida de Artefactos}: Se guardan los modelos entrenados \texttt{rf\_stacking.joblib}, \texttt{xgb\_stacking.joblib}, \texttt{lgbm\_stacking.joblib} y \texttt{meta\_learner.joblib} en el directorio \texttt{models/}.
\end{enumerate}

\subsection{Anexo D: Procedimiento de Evaluación, Validación 5-Fold y Gráficos}
\begin{enumerate}[leftmargin=*]
    \item \textbf{Generación de la Tabla Completa y Gráficos Senior (300 DPI)}: Se ejecutan los scripts de evaluación comparativa:
    \begin{lstlisting}[language=bash]
    python generate_senior_plots.py
    python build_complete_5phase_project.py
    \end{lstlisting}
    \item \textbf{Resultados Guardados}: Todas las imágenes en alta resolución y archivos CSV tabulares de resumen se guardan automáticamente en la carpeta \texttt{results/reports\_and\_plots/}.
\end{enumerate}

\subsection{Anexo E: Especificaciones del Entorno de Hardware y Software}
\begin{table}[H]
\centering
\caption{Especificaciones Técnicas del Entorno de Experimentación}
\label{tab:especificaciones_hardware}
\small
\begin{tabular}{lp{9.5cm}}
\toprule
\textbf{Componente} & \textbf{Especificación / Versión} \\
\midrule
Procesador (CPU) & Intel Core i7 / AMD Ryzen 7 de 12a Generación \\
Memoria RAM & 32 GB DDR5 4800 MHz \\
Tarjeta Gráfica (GPU) & NVIDIA GeForce RTX 5070 Ti Laptop GPU (CUDA 12.1) \\
Sistema Operativo & Windows 11 Home / Pro (64 bits) \\
Entorno Python & Miniconda3 - Entorno Virtual \texttt{investigacion} (Python 3.10) \\
Librerías Clave & PyTorch 2.2, ONNXRuntime 1.17, LightGBM 4.3, XGBoost 2.0, Scikit-Learn 1.4, SHAP 0.45, Streamlit 1.31, FastAPI 0.109 \\
\bottomrule
\end{tabular}
\end{table}

\subsection{Anexo F: Cuadernos de Jupyter Compilados (Notebooks de las 5 Fases)}
Los 5 cuadernos de Jupyter ejecutan el procedimiento metodológico completo y se ubican en la carpeta \texttt{notebooks/}:
\begin{enumerate}
    \item \texttt{01\_EDA\_Preprocessing\_and\_Feature\_Selection.ipynb}
    \item \texttt{02\_Model\_Architecture\_and\_Design.ipynb}
    \item \texttt{03\_Training\_Quantization\_and\_XAI.ipynb}
    \item \texttt{04\_Model\_Evaluation\_and\_Error\_Analysis.ipynb}
    \item \texttt{05\_State\_of\_the\_Art\_Comparison\_and\_Efficacy.ipynb}
\end{enumerate}
